% Generated draft: requires mathematical and source review.
\documentclass[11pt]{article}
\usepackage[margin=25mm]{geometry}
\usepackage{fontspec}
\setmainfont{Latin Modern Roman}
\usepackage{amsmath,amssymb,mathtools,unicode-math}
\setmathfont{Latin Modern Math}
\usepackage{enumitem,needspace,xurl,hyperref}
\setlength{\parindent}{0pt}
\setlength{\parskip}{5pt}
\setlength{\emergencystretch}{3em}
\newcommand{\N}{\mathbb{N}}
\newcommand{\Z}{\mathbb{Z}}
\newcommand{\Q}{\mathbb{Q}}
\newcommand{\R}{\mathbb{R}}
\newcommand{\C}{\mathbb{C}}
\newcommand{\F}{\mathbb{F}}
\newcommand{\A}{\mathbb{A}}
\newcommand{\PP}{\mathbb{P}}
\newcommand{\E}{\mathbb{E}}
\newcommand{\eps}{\varepsilon}
\newcommand{\dd}{\,\mathrm{d}}
\begin{document}
\section*{UnsolvedMath ID: 20003320}
\textbf{Source problem code:} catalogue entry
\subsection{Definitions and mathematical statement}
This record is historical context rather than the statement of Conjecture 34. The aid identifies its first sentences with a preceding positive-definite shrinking-gap question, and its following sentences with the distinct indefinite Oppenheim theorem. The source's ``indefinite definite'' wording is contradictory and is not adopted as a mathematical hypothesis.

The supplied special-case theorem is precise. Let \(P\) be an integer-valued positive-definite quadratic form on \(\Z^r\). Suppose there are integers \(M\geq1\), a residue \(a\pmod M\), and \(N_0\) such that every integer \(N\geq N_0\) with \(N\equiv a\pmod M\) is represented by \(P\). For every irrational \(\alpha>0\), set
\[
 Q(y,n)=P(y)+\alpha n^2,\qquad(y,n)\in\Z^r\times\Z.
\]
Since \(Q\) is positive definite, its distinct integral values form a locally finite increasing sequence \(v_0<v_1<\cdots\) tending to infinity. The aid asserts
\[
 v_{j+1}-v_j\longrightarrow0.
\]
This concerns consecutive \emph{distinct} values, not multiplicities. Equivalently, sufficiently far out, every interval of any fixed positive length meets the value set. No Diophantine condition on \(\alpha\) beyond positivity and irrationality is imposed.

For the supplied four-dimensional family, take \(P(y)=y_1^2+y_2^2+y_3^2\). Legendre's three-square theorem represents every positive integer congruent to \(1\pmod8\), so the progression hypothesis holds and
\[
 Q(x)=x_1^2+x_2^2+x_3^2+\alpha x_4^2
\]
has shrinking gaps for every irrational \(\alpha>0\). This does not establish the corresponding assertion for every positive-definite quaternary form.

The original source's historical status claim is that the preceding Conjecture 34 was proved in dimensions \(d\geq9\) in [15], then extended by that method to \(d\geq5\), while \(d=3,4\) were still open. The full preceding conjecture is not reproduced, and these are source-era assessments, not an independent current-status claim.

Separately, the Oppenheim discussion is read in its intended indefinite sense: if \(Q\) is a nondegenerate indefinite real quadratic form in dimension \(d\geq3\) and is not a nonzero real multiple of a rational form, the source attributes to Margulis [129] the density
\[
 \{Q(x):x\in\Z^d\}\text{ is dense in }\R.
\]
It says that proof was not effective. That indefinite density theorem is not the positive-definite shrinking-gap assertion and is not presented as a solution of the missing Conjecture 34. The supplied progression result is only a special family within that latter program.
\subsection{Short English statement}
The source mixes historical shrinking-gap status with a separate indefinite theorem. The supplied progression criterion proves shrinking gaps for one explicit four-dimensional family for every positive irrational coefficient. It does not solve all quaternary forms, and indefinite density is kept separate.
\subsection{Sources}
\begin{itemize}
\item Local UnsolvedMath registry, permanent ID \(20003320\), exact source code \texttt{catalogue entry}; canonical \texttt{aim-other-notes.json}, \texttt{notes[126]}, AIM domain Other, registry category Miscellaneous, source tag \texttt{conjecture}.
\item AIM workshop ``Emerging applications of measure rigidity,'' assigned item 34 context: \url{https://aimath.org/WWN/measrigid/measrigid.pdf}. Statement status \texttt{reconstructed\_unverified}; the record contains historical prose rather than the preceding conjecture. No PDF was fetched.
\item Source-local [15] concerns the historical positive-definite dimensional result; [129] is attributed to Margulis for Oppenheim. The aid invokes Legendre's three-square theorem for its special case. Full authoritative bibliographic entries are not supplied. The original literature field is empty.
\item Research attempt 1: result \texttt{context\_only}, problem status at run \texttt{not\_a\_problem}, rigor label \texttt{proved}, novelty confidence \texttt{low}; classification \texttt{OPEN-TRIAGE}, registry status \texttt{open}, suggested difficulty L2. Created and updated 14 August 2026. These are supplied assessments, not a resolution of every dimension-four form.
\item Attribution: Ulam AI / UnsolvedMath Contributors, CC BY 4.0. The machine-generated research aid is not peer reviewed and requires independent verification before citation.
\end{itemize}
\end{document}
